% \documentclass[12pt,oneside]{amsart}
\documentclass{scrartcl}

%\documentclass{article}

\usepackage{amsthm,amsmath,amssymb}
% \usepackage[numbers]{natbib}

%\usepackage[colorlinks]{hyperref}
\usepackage{hyperref}
\hypersetup{
    colorlinks,%
    citecolor=black,%
    filecolor=black,%
    linkcolor=black,%
    urlcolor=black
}
\usepackage{hypernat}
\usepackage[top=2.5cm, bottom=2.5cm, left=2.8cm,
right=2.8cm]{geometry}
\usepackage{graphicx} 
\usepackage{pdfpages}

%\setlength{\parskip}{0pt}
%\setlength{\parsep}{0pt}
%\setlength{\headsep}{0pt}
%\setlength{\topskip}{0pt}
%\setlength{\topmargin}{0pt}
%\setlength{\topsep}{0pt}
%\setlength{\partopsep}{0pt}

\linespread{1}

%\usepackage{marvosym}

%\textwidth = 6.5 in \textheight = 9.2 in \oddsidemargin = 0.0 in
%\evensidemargin = 0.0 in \topmargin = -0.0 in \headheight = 0.0 in
%\headsep = 0.0in \parskip = 0.0in \parindent = 0.0in
%\def\baselinestretch{0.871}


\newtheorem{theorem}{Theorem}[section]
\newtheorem{proposition}[theorem]{Proposition}
\newtheorem{problem}[theorem]{Problem}
\newtheorem{fact}[theorem]{Fact}
\newtheorem{lemma}[theorem]{Lemma}
\newtheorem{defn}[theorem]{Definition}
\newtheorem{corollary}[theorem]{Corollary}
\newtheorem{remark}{Remark}[section]
\newtheorem{example}[theorem]{Example}

\newcommand{\E}{{\mathbb E}}
\newcommand{\se}{{\mathcal{E}}}
\newcommand{\R}{{\mathbb R}}
\renewcommand{\P}{{\mathbb P}}
\newcommand{\ac}{{\mathcal{A}}}
\newcommand{\A}{{\mathcal{A}}}
\newcommand{\B}{{\mathcal{B}}}
\newcommand{\C}{{\mathcal{C}}}
\newcommand{\M}{{\mathcal{M}}}
\newcommand{\Mf}{{\mathfrak{M}}}
\newcommand{\T}{{\mathcal{T}}}
\newcommand{\Z}{{\mathcal{Z}}}
\newcommand{\K}{{\mathcal{K}}}
\newcommand{\lc}{{\mathcal{L}}}
\newcommand{\Ps}{{\mathcal{P}}}
\newcommand{\G}{{\mathcal{G}}}
\newcommand{\pa}{{\mathring{p}}}
\newcommand{\F}{{\cal F}}
\newcommand{\samp}{{\mathcal{X}}}
\newcommand{\X}{{\mathcal{X}}}
\newcommand{\hi}{{\hat{\Phi}}}
\newcommand{\data}{{\text{data}}}

\newcounter{rcnt}[section]
\renewcommand{\thercnt}{(\roman{rcnt})}

\def\qt#1{\qquad\text{#1}}

\def\argmin{\mathop{\rm argmin}}
\def\argmax{\mathop{\rm argmax}}


\setlength{\parskip}{1.4 \medskipamount}
%\sloppy
%\linespread{1.3}

\begin{document}

\title{Course List} 
% \subtitle{Fall 2026, UC Berkeley}
\author{Boting Zhang}
\date{}
\maketitle

\textbf{Course Taken:}
\begin{enumerate}
    \item \textbf{MAT 102 \& MAT 103 Calculus I, II}
    \\Main Contents: Limits, Derivatives, Integrals, Series, Multivariable Calculus, Vector Calculus.
    \\Reference: James Stewart, \textit{Calculus: Early Transcendentals}, 8th Edition, Cengage Learning, 2015.
    \item \textbf{MAT 104 \& MAT 105 Linear Algebra I, II}
    \\Main Contents: Vector Spaces, Linear Transformations, Eigenvalues and Eigenvectors, 
    Inner Product Spaces, Diagonalization and Quadratic Forms, Jordan Canonical Form, Numerical Methods.
    \\Reference: Howard Anton, \textit{Elementary Linear Algebra}, 11th Edition, Wiley, 2013.
    \item \textbf{MAT 107 Probability Theory}
    \\Main Contents: Probability Spaces, Random Variables, Expectation and Variance, Continuous Random Variables, 
    Joint Distributions, Limit Theorems.
    \\Reference: Sheldon Ross, \textit{A First Course in Probability}, 10th Edition, Pearson, 2018.
    \item \textbf{MAT 109 Ordinary Differential Equations (ODEs)}
    \\Main Contents: First-Order ODEs, Second-Order ODEs, Systems of ODEs, Laplace Transforms, Series Solutions,
    \\Reference: Boyce, DiPrima, \textit{Elementary Differential Equations and Boundary Value Problems}, 10th Edition, Wiley, 2012.
    \item \textbf{MAT 106 Discrete Mathematics}
    \\Main Contents: Logic and Proofs, Set Theory, Counting, Relations, Number Theory, Algorithms.
    \\Reference: Kenneth H. Rosen, \textit{Discrete Mathematics and Its Applications}, 7th Edition, McGraw-Hill, 2011.
    \item \textbf{MAT 101 Algebra and Analytic Geometry}
    \\Main Contents: Number systems, polynomials, sequences and series, partial 
    fractions, complex numbers, polar coordinates, parametric curves, conics, 
    geometry of three space.
    \\Reference: Julie Miller and Donna Gerken, \textit{College Algebra}, McGraw-Hill, 2023.
    \\James Stewart, Daniel K. Clegg and Saleem Watson, \textit{Calculus: Early Transcendentals}, Cengage Learning, 2020. 
    \item \textbf{BSC 135 Python Programming}
\end{enumerate}
\newpage


\textbf{Course Taking:}
\begin{enumerate}
    \item \textbf{MATH 126 Introduction to Partial Differential Equations (PDEs)}
    \\Main Contents: First-Order PDEs, Wave, Heat and Laplace equations, Fourier Series.
    \\Reference: 
    Borthwick, \textit{Introduction to Partial Differential Equations}, 2nd Edition, Springer, 2016.
    \\
    Lawrence C. Evans, \textit{Partial Differential Equations}, 2nd Edition, American Mathematical Society, 2010.
    \item \textbf{MATH 118 Fourier Analysis, Wavelets, and Signal Processing}
    \\Main Contents: Fourier Series, Fourier Transform, Wavelet Transform, Signal Processing Applications.
    \\Reference: Boggess and Narcowich, \textit{A First Course in Wavelets with Fourier Analysis}, 2nd Edition, Wiley, 2009.
    \item \textbf{STAT 248 Analysis of Time Series}
    \\Main Contents: Regressions, Spectral Analysis, ARIMA Models, Multivariate Time Series, Neural Networks.
    \\Reference: Shumway and Stoffer, \textit{Time Series Analysis and Its Applications: With R Examples}, 4th Edition, Springer, 2017.
\end{enumerate}


\newpage
\textbf{Course Want to take:}
\\2027 Spring
\begin{enumerate}
    \item \textbf{Mathematical Analysis I}
    \\Reference: Walter Rudin, \textit{Principles of Mathematical Analysis}, 3rd Edition.
    \item \textbf{Abstract Algebra}
    \\Reference: Michael Artin, \textit{Algebra}, 2nd Edition, Pearson, 2011.
    \item \textbf{Statistics}
    \item \textbf{Major Electives(ME): C++} 
\end{enumerate}


2027 Fall:
\begin{enumerate}
    \item \textbf{Mathematical Analysis II}
    \\Reference: Walter Rudin, \textit{Principles of Mathematical Analysis}, 3rd Edition.
    \item \textbf{Complex Analysis}
    \item \textbf{ME: Advanced Linear Algebra}
    \\Main Contents:  Abstract
vector spaces, linear transformations, the algebra of polynomials,
determinants and permutations, elementary canonical forms,
invariant subspaces, simultaneous triangulation, simultaneous
diagonalization, rational form, Jordan form, linear functionals and
adjoints, unitary operators, normal operators, operators on inner
product spaces, spectral theory, bilinear forms.
    \item \textbf{ME: Data Structures} \label{sec:ME:Data Structures}
\end{enumerate}




2028 Spring:
\begin{enumerate}
    \item \textbf{ME: Real Analysis}
    \\Main Contents: Functions of bounded variations, Riemann-Stieltjes
integrals, Lebesgue measure and outer measure on $\R^n$, Lebesgue
measurable sets, Lebesgue measurable functions, convergence in
measure, Lebesgue integrals, relations between Lebesgue integral
and Riemann-Stieltjes integral, Fubini's theorem, Tonelli's theorem,
Lebesgue's differentiation theorem, Vitali covering lemma, $L^p$
classes, Holder's inequality, Minkowski's inequality, Banach space,
metric space, Hilbert space, approximations of the identity,
Hardy-Littlewood maximal function
    \item \textbf{ME: Algebraic Topology}
    \\Main Contents: Fundamental
group, van Kampen theorem, covering spaces, deck
transformations, simplicial homology, singular homology, exact
sequences, Mayer-Vietoris sequences, cohomology, K¨unneth
formula, Poincar´e duality, homotopy group, Whitehead’s theorem,
fibre bundles.
    \item \textbf{ME: Design and Analysis of Algorithms}
    \\Main Contents: Introduction to algorithms, asymptotic analysis,
Analyzing Algorithms, probabilistic analysis, amortized analysis,
empirical analysis, recursive algorithm and recursive equation,
divide-and-conquer algorithms, dynamic programming, greedy
algorithms, graph algorithms, and NP-complete theory,
Backtracking, Branch and Bound.
    \item \textbf{ME: Differential Manifolds}
    \\Main Contents: Definitions and examples of manifolds, tangent spaces, partitions of
unity, tangent bundles, vector fields, Lie brackets, differential
forms, orientability, manifolds with boundary, integration on
manifolds, Stokes' theorem on manifolds, and a brief introduction
to Lie groups and Lie algebras.
    \item \textbf{Final Year Project}
\end{enumerate}

2028 Fall:
\begin{enumerate}
    \item \textbf{ME: Mechanics}
    \item \textbf{ME: Numerical Analysis} or \textbf{ME: Stochastic Processes}
    \\Main Contents(Numerical Analysis): Newton-Raphson method, zeros of polynomials,
interpolation by Lagrange polynomials, Hermite interpolation,
cubic spline interpolation, numerical differentiation, numerical
integration, Runge-Kutta methods, stability, solving linear systems,
approximation theory, shooting method, finite-difference method.
    \\Main Contents(Stochastic Processes):Markov chains, exponential distributions,
Poisson processes, continuous time Markov chains, renewal theory,
queueing theory, reliability theory, Brownian motions, stationary
processes.
    \item \textbf{Final Year Project}
\end{enumerate}    
    


\newpage











\end{document}

%%% Local Variables:
%%% mode: latex
%%% TeX-master: t
%%% End:
